% ============================================================
\chapter{Conclusion}
\label{ch:conclusion}
% ============================================================

This report presents the design and optimisation of a resistive rear window defroster using coupled multiphysics finite element modelling in COMSOL Multiphysics~6.2. The design was developed iteratively through six simulations addressing thermal performance, material selection, structural integrity, transient response, trace pattern optimisation, and weather condition robustness.

The key findings of this study are:

\begin{enumerate}
    \item \textbf{The number of heating traces is the most influential design parameter.} Increasing the trace count from 12 to 20 raises $T_\text{min}$ by over \SI{6}{\kelvin} at \SI{12}{\volt}, whereas increasing voltage from 12 to \SI{14.5}{\volt} at 14~traces yields only \SI{1.2}{\kelvin}.

    \item \textbf{Pd--Ag alloy ($\sigma = \SI{2.3e6}{\siemens\per\metre}$) is the optimal trace material.} It reaches the dew point with 14~traces at \SI{12}{\volt}, while Nichrome~V requires 18--20~traces. Bulk conductivity materials (Silver, Gold) overheat catastrophically and are not viable for this geometry.

    \item \textbf{Both shortlisted materials pass the stress analysis} with safety factors exceeding $3\times$, confirming that thermal expansion mismatch between the traces and glass does not pose a structural risk.

    \item \textbf{The recommended design} is:
    \begin{itemize}
        \item Material: Pd--Ag alloy
        \item Number of traces: 20
        \item Voltage: \SI{12}{\volt} (nominal battery)
        \item Steady-state $T_\text{min}$: \SI{288.2}{\kelvin} (\SI{5}{\kelvin} above dew point)
        \item Defrost time: 12~minutes
        \item Stress safety factor: $\sim 3.1 \times$
    \end{itemize}

    \item \textbf{Alternative high-performance design:} Pd--Ag alloy, 22~traces, \SI{14.5}{\volt} --- achieves 8-minute defrost time when the engine is running.

    \item \textbf{The final design is robust across all tested weather conditions (Sim~6).} Five scenarios ranging from mild autumn (\SI{10}{\celsius}) to extreme cold (\SI{-20}{\celsius}) and freezing rain ($h_\text{ext} = \SI{40}{\watt\per\metre\squared\per\kelvin}$) were tested, each with a physically realistic dew point \cite{lawrence2005}. All five scenarios defrost within \SI{8}{\minute} with a minimum steady-state margin of \SI{+5.4}{\kelvin} above the dew point, confirming that the design is not tuned to a single favourable condition.

    \item \textbf{Chebyshev trace spacing halves the defrost time.} The non-uniform Chebyshev spacing (\Cref{eq:chebyshev}) with exponent $p = 1.3$ reduces the defrost time from \SI{12}{\minute} to \SI{6}{\minute} and raises $T_\text{min}$ by \SI{6.8}{\kelvin} --- without changing the number of traces, material, or voltage. The discrete zoned approach (6/8/6) provides only a modest \SI{1}{\minute} improvement, confirming that a continuous optimisation is needed. The optimal exponent ($p = 1.3$) was found via a 16-point parametric sweep, exploiting COMSOL's ability to parametrically update the geometry.

    \item \textbf{Final recommended design:}
    \begin{itemize}
        \item Material: Pd--Ag alloy ($\sigma = \SI{2.3e6}{\siemens\per\metre}$)
        \item Number of traces: 20
        \item Spacing: Chebyshev with $p = 1.3$
        \item Voltage: \SI{12}{\volt} (nominal battery)
        \item Steady-state $T_\text{min}$: \SI{295.0}{\kelvin} (\SI{12}{\kelvin} above dew point)
        \item Defrost time: \SI{6}{\minute}
        \item Stress safety factor: $\sim 3.1 \times$
        \item Window obscuration: 4.25\% ($<5\%$ SAE~J216 limit)
    \end{itemize}
\end{enumerate}

The simulation methodology --- coupling Electric Currents, Heat Transfer, and Solid Mechanics with parametric sweeps over material properties, geometry, and operating conditions --- provides a systematic framework for defroster design. The model was validated through mesh convergence, back-of-envelope analytical estimates (agreeing to within 0.03\%), and comparison with published experimental data.

The key insight of this study is that \textbf{geometry optimisation} (trace placement) can be equally or more impactful than material or electrical optimisation. The Chebyshev spacing approach provides a mathematically principled, single-parameter design specification that can be applied to any rectangular heated-glass application. Future work should extend the model to 3D curved glass, validate against physical prototypes, and investigate the effect of varying interior convection ($h_\text{int}$) with different cabin HVAC states.
